\subsection{The modular quotients after pullback to the Legendre line}

\begin{statement}{A.53}{Proved}{Exact level-three, four, five and seven maps}
Put $k=\bar\Q$.  The projective geometric quotients of the Legendre
monodromy at levels $N=3,4,5,7$, and the smooth compactifications of
their connected covers of $U_k$, are as follows:
\[
\begin{array}{c|c|c|c|c}
N&\text{linear image}&\text{projective image }Q_N
 &\text{compactification}&\text{genus}\\ \hline
3&\operatorname{SL}_2(\mathbf F_3)&A_4&X(6)_k&1\\
4& (\Z/2)^2&(\Z/2)^2&X(4)_k&0\\
5&\operatorname{SL}_2(\mathbf F_5)&A_5&X(10)_k&13\\
7&\operatorname{SL}_2(\mathbf F_7)&
       \operatorname{PSL}_2(\mathbf F_7)&X(14)_k&49.
\end{array}
\tag{A.53.1}
\]
Here $X(N)$ has its usual symplectic modular meaning over a field
containing the relevant roots of unity.  The geometric assertions in
the table do not assert that every arithmetic twist has a rational
point or is isomorphic over $\Q$ to that chosen modular model.

The group $S_4\simeq\operatorname{SL}_2(\Z/4)/\{\pm I\}$ belongs to
$X(4)\to X(1)$ over the $j$-line.  The component obtained after
pullback along $X(2)\to X(1)$ has group
\[
 \ker\bigl(S_4\longrightarrow S_3\bigr)=(\Z/2)^2.
\tag{A.53.2}
\]
At odd $N$ in the table the connected pullback is $X(2N)$.
In particular, the genus-zero curves $X(3),X(5)$ and the genus-three
curve $X(7)$ in L10 are the curves before this pullback.

For each $N$ let
\[
 N_N=\ker\bigl(\Pi\xrightarrow{\rho_{\rm Leg,N}}
      \operatorname{GL}_2(\Z/N)\longrightarrow
      \operatorname{PGL}_2(\Z/N)\bigr).
\tag{A.53.3}
\]
This is a closed normal subgroup of $E_U$, and the section in the
question has the actual quotient section
\[
       G_\Q\xrightarrow{s}E_U\longrightarrow E_U/N_N,
       \qquad \Pi/N_N=Q_N.
\tag{A.53.4}
\]
Thus A.51 applies with precisely the four groups in (A.53.1).
For $N=5,7$ these groups are perfect.  For $N=4$ the quotient
already factors through the full geometric abelianization; adjoining
it to that abelianization adds no further quotient information.
\end{statement}
\begin{proof}
We first give the maps on the base curves, including the loci on
which they are etale.  Set
\[
 J(t)=256\frac{(t^2-t+1)^3}{t^2(1-t)^2},\qquad
 W=\mathbb P^1_j\setminus\{0,1728,\infty\},\qquad
 U^\circ=J^{-1}(W).
\tag{A.53.5}
\]
Thus $U^\circ$ is obtained from $U$ by deleting the three harmonic
points and the degree-two equianharmonic point.  The identity
\[
 J(t)-1728=
 64\frac{(t-2)^2(t+1)^2(2t-1)^2}{t^2(1-t)^2}
\tag{A.53.6}
\]
shows the ramification over $1728$; the numerator of $J$ shows the
ramification over zero.  Its poles have order two at $0,1,\infty$.
The six transformations
\[
 t,\quad 1-t,\quad t^{-1},\quad (1-t)^{-1},\quad
 (t-1)/t,\quad t/(t-1)
\tag{A.53.7}
\]
preserve $J$.  They give all automorphisms of its degree-six function
field extension.  This proves that $U^\circ\to W$ is a finite etale
$S_3$-cover, and that its compactification is $X(2)\to X(1)$ with the
stated coordinate.  On an elliptic curve the six transformations
are obtained by permuting the three nonzero points of order two and
then sending the selected first two roots to zero and one.  This
also identifies the map with the forgetful modular map, rather than
only identifying two abstract degree-six covers.

For the modular comparison one may work over $k$ throughout.
The definition of a symplectic level structure and its modular
curve is recalled in \cite[Section 8, pp.~99--100]{LegMilneMF2017}.
Changing a symplectic frame acts through
$\operatorname{SL}_2(\Z/N)$.  The two frames $\alpha$ and $-\alpha$
define isomorphic pairs, by the actual elliptic-curve isomorphism
$[-1]$.  At a generic $j$ the elliptic curve has precisely these
two automorphisms.  Consequently the generic deck group over the
$j$-line is
\[
 P_N=\operatorname{SL}_2(\Z/N)/\{\pm I\}.
\tag{A.53.8}
\]
Removing $j=0,1728,\infty$ gives its finite etale locus.

At level four reduction modulo two gives the exact sequence
\[
 1\longrightarrow V\longrightarrow P_4
   \longrightarrow \operatorname{SL}_2(\mathbf F_2)
   \longrightarrow1,
 \qquad
 V=\left\{\left[\begin{pmatrix}1&2a\\2b&1\end{pmatrix}\right]
                 :a,b\in\mathbf F_2\right\}.
\tag{A.53.9}
\]
The multiplication on $V$ is $(a,b)+(a',b')$, so $V$ is the Klein
four group.  Before quotienting by $\{\pm I\}$ the reduction
kernel consists of $I+2M$ with $\operatorname{tr}(M)=0$ in
$\mathbf F_2$.  It has eight elements; quotienting identifies the
two diagonal choices and gives exactly the four matrices in
(A.53.9).

For completeness the identification $P_4\simeq S_4$ can be made
without choosing a name for an abstract group of order 24.
Let $P_4$ act by conjugation on the four cyclic subgroups generated,
respectively, by the projective classes of
\[
 \begin{pmatrix}0&1\\3&1\end{pmatrix},\quad
 \begin{pmatrix}0&1\\3&3\end{pmatrix},\quad
 \begin{pmatrix}1&1\\1&2\end{pmatrix},\quad
 \begin{pmatrix}1&3\\3&2\end{pmatrix}.
\tag{A.53.10}
\]
These are its four order-three subgroups.  Conjugation by
$S=\left(\begin{smallmatrix}0&-1\\1&0\end{smallmatrix}\right)$
acts as $(12)$, and conjugation by
$T=\left(\begin{smallmatrix}1&1\\0&1\end{smallmatrix}\right)$
acts as $(1342)$, in the displayed ordering.  These permutations
generate $S_4$, while $P_4$ has order 24.  The homomorphism is
therefore an isomorphism.  Its reduction quotient is $S_3$, and
its kernel is (A.53.9), proving the exact meaning of (A.53.2).

Let $X(N)^\circ$ be the inverse image of $W$.  The component of
\[
                 X(4)^\circ\times_W U^\circ
\tag{A.53.11}
\]
selected by the forgetful map $X(4)\to X(2)$ is the graph of that
map.  It is isomorphic to $X(4)^\circ$, and its projection to
$U^\circ$ has deck group $V$.  Over $k$ the entire fibre product
is the disjoint union of six such components, indexed by the six
automorphisms in (A.53.7).  This is the usual fibre-product
decomposition of a Galois cover along its intermediate normal
quotient, applied here to the explicit reduction map (A.53.9).

The selected cover over all of $U$ has an elementary model:
\[
 C:\ x^2+y^2=1,\qquad t=x^2,
\tag{A.53.12}
\]
with the points over $t=0,1,\infty$ removed.  Independent changes
of the signs of $x,y$ give its group $V$.  Indeed the level-four
Legendre matrices from A.16, restricted to the geometric group,
are
\[
 \rho_{\rm Leg,4}(g)=
 \begin{pmatrix}
 1&2\kappa_2(1-t)(g)\\2\kappa_2(t)(g)&1
 \end{pmatrix},\qquad g\in\Pi.
\tag{A.53.13}
\]
Their common kernel fixes $k(t)(\sqrt t,\sqrt{1-t})$, the function
field of (A.53.12).  This supplies the comparison with the actual
Legendre representation.  The parametrization and composed map
are
\[
 \begin{gathered}
 x=\frac{2z}{1+z^2},\qquad
 y=\frac{1-z^2}{1+z^2},\qquad
 t=\frac{4z^2}{(1+z^2)^2},\\
 J(t(z))=16\frac{(z^8+14z^4+1)^3}{z^4(z^4-1)^4}.
 \end{gathered}
\tag{A.53.14}
\]
This gives a rational genus-zero model and the actual degree-four
and degree-24 maps in the level-four diagram.

Now let $N$ be odd.  The Chinese remainder isomorphism gives
\[
 \operatorname{SL}_2(\Z/2N)/\{\pm I\}
  \xrightarrow{\ \sim\ }
 \operatorname{SL}_2(\mathbf F_2)
     \times\bigl(\operatorname{SL}_2(\Z/N)/\{\pm I\}\bigr).
\tag{A.53.15}
\]
Here $-I$ reduces to $I$ modulo two; this proves the asserted
quotient formula.  In terms of the level structures, the maps are
induced by
\[
 E[2N]\longrightarrow E[2]\oplus E[N],\qquad
                   P\longmapsto([N]P,[2]P).
\tag{A.53.16}
\]
If $Na+2b=1$, the inverse on the direct sum is
$(P_2,P_N)\mapsto aP_2+bP_N$; the restrictions of the Weil
pairing give the corresponding symplectic frames.  Thus the two
projections of (A.53.15) are the actual forgetful maps of curves.
They show that
\[
 X(2N)^\circ\xrightarrow{\ \sim\ }
                     X(N)^\circ\times_W U^\circ
\tag{A.53.17}
\]
over $k$.  In particular this fibre product is connected, and its
deck group over $U^\circ$ is $P_N$.

One can also see the same comparison directly on the complex
uniformizing subgroups:
\[
 \overline{\Gamma(2)}\cap\overline{\Gamma(N)}
                 =\overline{\Gamma(2N)}
       \quad\text{in }\operatorname{PSL}_2(\Z).
\tag{A.53.18}
\]
The bar denotes the image after dividing by the central sign.
For the reverse inclusion in (A.53.18), choose the sign of a
representative which is the identity modulo $N$; changing its
sign does not alter its identity reduction modulo two.  It is
then the identity modulo $2N$ by the Chinese remainder theorem.
This argument proves the equality of subgroups, including its
central sign.

A.17 gives the geometric Legendre generators $U(2),L(-2)$.
For $N=3,5,7$, their reductions generate
$\operatorname{SL}_2(\mathbf F_N)$: their powers give all upper
and lower elementary matrices, and row reduction expresses every
determinant-one matrix as a product of these.  Projectivizing
therefore gives $P_N$.  At $N=3$ its faithful action on
$\mathbb P^1(\mathbf F_3)$ consists of even permutations and has
order 12, giving $P_3=A_4$.

An explicit identification $P_5=A_5$ is conjugation on its five
Klein four subgroups.  For example start with the subgroup
generated by the projective classes of
$\left(\begin{smallmatrix}0&1\\4&0\end{smallmatrix}\right)$ and
$\left(\begin{smallmatrix}0&2\\2&0\end{smallmatrix}\right)$,
and order its five conjugates by successive conjugation by
$U(2)$.  The permutations induced by $U(2),L(-2)$ are
$(12345)$ and $(14523)$.  They generate $A_5$: their quotient
contains the three-cycle $(142)$, and their subgroup has order
divisible by both three and five.  A subgroup of $A_5$ of order
15 would have a normal subgroup of order five and hence commute
with its order-three subgroup, giving an impossible permutation
of order 15 on five letters.  An order-30 subgroup would give
a nontrivial homomorphism $A_5\to\Z/2$, whereas $A_5$ is generated
by three-cycles.  Its order is therefore 60.  Since $P_5$ also
has order 60 this action is faithful.  The finite matrix
calculations, including the list of all five subgroups, are
retained in the certificate named below.

At these odd levels the central matrix $-I$ belongs to the linear
geometric image and gives the exact sequence
\[
 1\longrightarrow\{\pm I\}\longrightarrow
 \operatorname{SL}_2(\mathbf F_N)\longrightarrow Q_N
 \longrightarrow1.
\tag{A.53.19}
\]
The linear frame cover and its projective quotient are related by
this central quotient.  At level four the actual image
(A.53.13) does not contain $-I$; its projective image has the
same four elements.  This accounts for the different central
behavior of the Legendre family at that level.

The finite projective covers just constructed extend over the
five deleted points of $U\setminus U^\circ$.  One verification is
to take the quotient of the finite etale frame scheme of
$\mathcal E[N]$ over $U$ by its central scalar action.  Its
geometric connected component is the cover computed above;
over $U^\circ$ it is identified by the projective frame with
(A.53.11) or (A.53.17).  Consequently normalization gives the
same cover over $U$.  This explains why (A.53.4) uses a section
of $U$ directly: no lift of $s$ to $U^\circ$ is being presumed.

The three projective inertia generators at the $\lambda$-cusps
have order $N$ for odd $N$, and order two in the level-four
image.  For the infinity generator this uses the central sign:
it is projectively unipotent, although its odd-level linear
order is $2N$.  If $d=|Q_N|$ and $e$ is this common projective
inertia order, Riemann--Hurwitz gives
\[
                    2g-2=d\bigl(1-3/e\bigr).
\tag{A.53.20}
\]
The orders are $d=12,4,60,168$ and $e=3,2,5,7$, respectively.
This proves the genus column in (A.53.1).  It agrees with the
principal-level genus formula in
\cite[Theorem 2.22 and the examples on p.~38]{LegMilneMF2017}.
The modular identification of $X(7)$ with the Klein quartic is
the one used in \cite[Sections 3--4]{LegPoonenSchaeferStoll2007};
the map from the relevant cover here to that curve is the
degree-six forgetful map $X(14)\to X(7)$ in (A.53.16).

Finally, normality in (A.53.3) follows because conjugation by
the arithmetic monodromy preserves its geometric projective
kernel.  More explicitly, for odd $N=3,5,7$ one has isomorphisms
of profinite extensions of $G_\Q$,
\[
 E_U/\ker(\rho_{{\rm Leg},N}|_\Pi)
  \simeq
 \{(M,\sigma)\in\operatorname{GL}_2(\mathbf F_N)\times G_\Q:
                         \det M=\chi_N(\sigma)\},
\tag{A.53.21}
\]
\[
 E_U/N_N\simeq
 \{([M],\sigma)\in\operatorname{PGL}_2(\mathbf F_N)\times G_\Q:
           \det[M]=[\chi_N(\sigma)]
                  \text{ in }\mathbf F_N^\times/(\mathbf F_N^\times)^2\}.
\tag{A.53.22}
\]
The maps are $(\rho_{{\rm Leg},N},\operatorname{pr})$ and its
projectivization.  Their kernels are the displayed geometric
kernels.  They are onto because the geometric images are
$\operatorname{SL}_2$ and $\operatorname{PSL}_2$, respectively:
above each $\sigma$ one chosen lift supplies one matrix with the
correct determinant, and multiplying by the geometric image
supplies every other such matrix or projective matrix.  The
continuous bijections are topological isomorphisms by compactness.
Thus (A.53.21), (A.53.22), and (A.53.19) give the precise relation
between the linear and projective arithmetic data of a section.

At level four the analogous map is even simpler:
\[
 E_U/N_4\xrightarrow{\ \sim\ }
       (\mathbf F_2)^2\times G_\Q,
 \qquad
 g\longmapsto
   \bigl(\kappa_2(t)(g),\kappa_2(1-t)(g),\operatorname{pr}(g)\bigr).
\tag{A.53.23}
\]
The two Kummer characters have constant coefficient module
$\mu_2$, and hence are homomorphisms.  Their common geometric
kernel is (A.53.13), and their geometric image is all of
$(\mathbf F_2)^2$, proving the isomorphism as above.  The
cyclotomic diagonal $\operatorname{diag}(\chi_4,1)$ of A.16 is
already specified by the $G_\Q$ coordinate.  In particular
$\Gamma_2\Pi\subset N_4$ and the fibre product in A.51 for this
quotient is exactly $E_1$.

For $N=5,7$ perfectness can be checked within the displayed
matrices.  Put $D=\operatorname{diag}(2,2^{-1})$.  The commutator
$[D,U(a)]=U(3a)$ and its lower-triangular counterpart contain
every elementary matrix since three is invertible in both
fields.  Thus $\operatorname{SL}_2(\mathbf F_N)$, and consequently
its projective quotient, is perfect.  At $N=3$ the derived
subgroup of $A_4$ is $V$, with quotient $\Z/3$; at $N=4$ the
group is already abelian.  These are exactly the inputs required
for the applications of A.51.
\end{proof}

\noindent\emph{Retained verification.}
\texttt{code/legendre\_modular\_levels.py} enumerates the
four finite linear and projective images, implements the conjugation
actions (A.53.10) and the five Klein subgroups, verifies (A.53.14)
and all six transformations (A.53.7), and computes (A.53.20) with
exact rational arithmetic.  Its complete output is retained in
\texttt{code/legendre\_modular\_levels\_output.txt}.
